\section{Linear decay and transfer to expanding tori}\label{sec:linear}

Fix one sufficiently small admissible $a>0$ and a profile $Q=Ae^{i\phi}$
from Proposition~\ref{profile:existence} for which
Theorem~\ref{spectrum:classification} holds.  All constants in this section
may depend on this fixed profile and on $a$.  No uniformity as $a\downarrow0$
is needed here.  Put $P=|Q|^{1/a}$ and write the real-linear potential as
\begin{equation}\label{lin:potential}
 \mathcal Vv=P\left(1+\frac1{2a}\right)v
       +\frac{P}{2a}e^{2i\phi}\overline v,
 \qquad
 \mathcal A=i\Delta-\frac12y\cdot\nabla-a+ib-i\mathcal V.
\end{equation}
The profile estimates give, for every nonnegative integer $j$,
\begin{equation}\label{lin:symbol-hypotheses}
 |\nabla^j Q(y)|\le C_j\langle y\rangle^{-2a-j},\qquad
 |\nabla^j\mathcal V(y)|\le C_j\langle y\rangle^{-2-j}.
\end{equation}
The second inequality applies to each of the two coefficients in
\eqref{lin:potential}.  We take $0<a<1$.  The argument below uses the
all-order estimates only up to a finite order, after that order has been
chosen.

\subsection{Spaces, products, and localization}

Set $s_0=6-a$.  For an integer $k>8$ to be fixed below, let $Y$ be the real
Hilbert space obtained by completing complex-valued Schwartz functions in
\begin{equation}\label{lin:Y}
 \|v\|_Y^2=\|v\|_{\dot H^{s_0}(\mathbb R^{12})}^2
                      +\|v\|_{\dot H^k(\mathbb R^{12})}^2.
\end{equation}
Homogeneous Fourier multipliers have their usual Plancherel normalization.
In particular, no global $L^2$ norm is included in \eqref{lin:Y}.
The choice $6-2a<s_0=6-a<6$ serves two purposes: the lower inequality
admits the profile tail of order $r^{-2a}$, while the upper inequality
makes the inverse Fourier weight integrable at low frequency. Under the
free similarity evolution this norm component will have squared-norm
rate $6-s_0-2a=-a<0$. The high order $k$ supplies the local regularity
and the remaining dissipation.

For $L\ge1$ put $\mathbb T_L^{12}=\mathbb R^{12}/(2\pi L\mathbb Z)^{12}$.
If $v(y)=\sum_{n\in\mathbb Z^{12}}v_ne^{in\cdot y/L}$, define
\begin{equation}\label{lin:YL}
 \|v\|_{Y_L}^2=(2\pi L)^{12}\sum_n
 \left[(L^{-2}+|n/L|^2)^{s_0}+|n/L|^{2k}\right]|v_n|^2.
\end{equation}
Thus $Y_L$ is $H^k(\mathbb T_L^{12})$ with a scale-dependent norm.
The regularization in its first summand is essential: a constant of value
$z$ has norm $(2\pi)^6L^a|z|$.
Throughout this section $\nabla^k$ is the \emph{ordered} derivative tensor:
its components are $\partial_{j_1}\cdots\partial_{j_k}$ for all ordered
$k$-tuples.  Consequently
$\|\nabla^kv\|_2^2=\|v\|_{\dot H^k}^2$ exactly.

\begin{lemma}[Uniform spatial estimates]\label{lin:spatial}
The embeddings $Y\hookrightarrow C_0(\mathbb R^{12})$ and
$Y_L\hookrightarrow L^\infty(\mathbb T_L^{12})$ are bounded, uniformly for
$L\ge1$.  Both spaces are algebras, with a uniform algebra constant for
$Y_L$.  For a ball or annulus of radius $R\ge1$, contained in one torus
cell when appropriate and with $R\le C_0L$, one has
\begin{equation}\label{lin:local-derivatives}
 \|\nabla^jv\|_{L^2(\text{ball or annulus})}
 \le C R^{\max(s_0-j,0)}\|v\|,
 \qquad 0\le j\le k.
\end{equation}
Here and below an undecorated norm denotes the applicable $Y$ or $Y_L$
norm, and constants may depend on the fixed geometric constant $C_0$.

If $F$ satisfies $|\nabla^jF|\le C_j\langle y\rangle^{-\beta-j}$, then
a smooth dyadic piece supported where $|y|\asymp R$ satisfies
\begin{equation}\label{lin:symbol-norm}
 \|\psi_RF\|\le C\big(R^{6-\beta-s_0}+R^{6-\beta-k}\big),
\end{equation}
also on the tori for $R\le C_0L$.  In particular $Q\in Y$, and for every
fixed $\chi\in C_c^\infty(B_1)$ the cutoff profiles $\chi(y/L)Q(y)$
have uniformly bounded $Y_L$ norms.
\end{lemma}

\begin{proof}
For the continuous Fourier transform, Cauchy--Schwarz uses
\[
 \int_{|\xi|\le1}|\xi|^{-2s_0}\,d\xi<\infty,
 \qquad
 \int_{|\xi|\ge1}|\xi|^{-2k}\,d\xi<\infty.
\]
Thus $\|\widehat v\|_1\le C\|v\|_Y$.  This identifies the completion in
\eqref{lin:Y} with functions vanishing at infinity, rather than with
distributions modulo polynomials.  In the periodic case the corresponding
inverse-weight sum, after division by $(2\pi L)^{12}$, is uniformly
bounded.  Indeed its portion $|n|\le L$ is bounded by
\[
 C L^{-12+2s_0}\sum_{|n|\le L}\langle n\rangle^{-2s_0}\le C,
\]
and the portion $|n|>L$ is bounded using $k>6$.
The convolution inequality
$|\xi+\zeta|^s\le C_s(|\xi|^s+|\zeta|^s)$, and its version for
$(L^{-2}+|\xi|^2)^{s/2}$, now prove the algebra estimates by Young's
inequality.  The volume factors in \eqref{lin:YL} cancel against the
Fourier $\ell^1$ estimate just proved.

For $j<s_0$, split at frequency $R^{-1}$.  The low-frequency part has
$j$th derivative bounded pointwise by $CR^{-a-j}\|v\|$.  This follows
by the same inverse-weight integral or sum; it includes the periodic
zero mode.  Multiplication by the square root of the volume of the
ball gives $CR^{s_0-j}\|v\|$.  On higher frequencies, Plancherel gives
that same bound because
$|\xi|^j\le R^{s_0-j}|\xi|^{s_0}$ there.  For $s_0\le j\le k$,
the global derivative norm is bounded by the two defining Sobolev
norms.  This proves \eqref{lin:local-derivatives}.

For a localized symbol, the integer derivative norms of order $m$ are
$O(R^{6-\beta-m})$, with the same bounds for derivatives falling on
$\psi_R$.  Fourier interpolation gives the fractional order in
\eqref{lin:symbol-norm}.  The additional periodic term is bounded by
\[
 L^{-s_0}\|\psi_RF\|_2
 \le C(R/L)^{s_0}R^{6-\beta-s_0},
\]
since the first norm in \eqref{lin:YL} is uniformly equivalent to
$\|\cdot\|_{\dot H^{s_0}}+L^{-s_0}\|\cdot\|_2$.
For $\beta=2a$ the first exponent in \eqref{lin:symbol-norm} is $-a$;
both exponents are negative.  Summing dyadic pieces proves the last
assertions, and also convergence of smooth compactly supported
approximations to $Q$ in $Y$.
\end{proof}

Fix a real $\chi\in C_c^\infty(B_1)$ with $0\le\chi\le1$ and
$\chi=1$ on $B_{1/2}$.  Define
\begin{equation}\label{lin:localization}
 J_Lv(y)=\chi(y/L)v(y),
\end{equation}
using the central-cell representative on its support and extending by
zero to $\mathbb R^{12}$.  There is a constant $C_J$, independent of
$L\ge1$, such that
\begin{equation}\label{lin:J-bound}
 \|J_Lv\|_Y\le C_J\|v\|_{Y_L}.
\end{equation}
Here is a direct verification, useful because a fractional derivative
does not localize pointwise.  Up to fixed Fourier normalization constants,
\[
 \widehat{J_Lv}(\xi)=L^{12}\sum_nv_n\widehat\chi(L\xi-n).
\]
For $s=s_0,k$, set $\zeta=L\xi$ and use
$|\zeta|^s\le C_s\langle n\rangle^s\langle\zeta-n\rangle^s$.
The mixed discrete-to-continuous convolution map with kernel
$\langle\zeta\rangle^s|\widehat\chi(\zeta)|$ is bounded from
$\ell^2$ to $L^2$: Cauchy--Schwarz uses the uniformly bounded sum of
its translates, followed by integration of the kernel.  It follows that
\[
 \|J_Lv\|_{\dot H^s}^2
 \le C L^{12-2s}\sum_n\langle n\rangle^{2s}|v_n|^2.
\]
For $s=k$ the extra term $L^{-2k}\|v\|_2^2$ is controlled by the
regularized low-order norm since $L^{-k}\le L^{-s_0}$.  This proves
\eqref{lin:J-bound}.

\subsection{Evolution and dissipation with a local observation}

The free operator $\mathcal A_0=\mathcal A+i\mathcal V$ has the group
\begin{equation}\label{lin:free-flow}
 S_0(s)v(y)=e^{(-a+ib)s}
 \left[e^{i(1-e^{-s})\Delta}v\right](e^{-s/2}y),\qquad s\in\mathbb R.
\end{equation}
Fourier approximation proves strong continuity on $Y$.  The potential
is bounded on $Y$: since $p=1+1/a$ is an odd integer, its coefficients
are polynomials in $Q,\overline Q$, and Lemma~\ref{lin:spatial} applies.
The integral equation
\begin{equation}\label{lin:mild}
 S(s)v=S_0(s)v-i\int_0^s S_0(s-t)\mathcal V S(t)v\,dt
\end{equation}
therefore constructs a strongly continuous real-linear group on $Y$.
Successive approximation on short intervals proves existence, and
Gronwall's inequality proves uniqueness and boundedness on every compact
time interval.

For the torus evolution, put $L_s=Le^{s/2}$ and
\begin{equation}\label{lin:torus-coefficients}
 Q_s(y)=\chi(y/L_s)Q(y),\qquad
 \mathcal V_s=D\big(|U|^{1/a}U\big)\big|_{U=Q_s}.
\end{equation}
Both are defined cellwise and periodically.  In particular
$\mathcal V_s=\chi(y/L_s)^{p-1}\mathcal V$ on the central cell.
The equation is
\begin{equation}\label{lin:torus-evolution}
 (\partial_s+\tfrac12y\cdot\nabla)v
       =i\Delta v-a v+ibv-i\mathcal V_s v.
\end{equation}
It is interpreted on the expanding torus: the left side differentiates
at fixed $y/L_s$.  Formula~\eqref{lin:free-flow} and a nonautonomous
version of \eqref{lin:mild} apply.  The uniform product bounds give a
unique propagator with bounds depending on a fixed slab length, but
independent of the starting scale $L\ge1$.

Let $B$ bound the pointwise real operator norms of $\mathcal V$ and all
$\mathcal V_s$.  It is finite and independent of $k$; one may take
$B=(1+1/a)\|P\|_\infty$.  We now choose and fix an integer
\begin{equation}\label{lin:k-choice}
 k>\max\{8,\,6-2a+2B+1\}.
\end{equation}
Subsequent constants may depend on this $k$.

\begin{lemma}[Compact errors]\label{lin:compact-error}
For every $\varepsilon>0$ there are $R_\varepsilon,C_\varepsilon<\infty$
such that
\begin{equation}\label{lin:compact-bound}
 \|\mathcal V_*v\|_{s_0,*}
 +\|\nabla^k(\mathcal V_*v)-\mathcal V_*\nabla^kv\|_2
 \le\varepsilon\|v\|+C_\varepsilon\|v\|_{L^2(B_{R_\varepsilon})}.
\end{equation}
Here $*$ denotes either the whole-space potential and homogeneous
low-order norm, or a torus potential at its current scale and the
regularized low-order norm in \eqref{lin:YL}.  The constants are uniform
for all sufficiently large current torus scales.  The potential acts
componentwise on the ordered tensor in the second term.
\end{lemma}

\begin{proof}
All coefficient derivatives have the bounds in
\eqref{lin:symbol-hypotheses}, uniformly for the cutoff potentials.
On an annulus of radius $R$, a lower Leibniz term with $j<k$ derivatives
on $v$ has norm at most
\begin{equation}\label{lin:top-tail}
 C_kR^{-2-k+j+\max(s_0-j,0)}\|v\|.
\end{equation}
Every exponent is negative: it is $-2-k+s_0$ for $j<s_0$ and at most
$-3$ otherwise.  Dyadic summation therefore makes these tail terms
arbitrarily small.

For the fractional part, use a smooth dyadic partition and apply
\emph{global} Fourier interpolation to the actual localized product
$z_R=\psi_R\mathcal V_*v$.  Leibniz' rule, including derivatives of
$\psi_R$, gives
\[
 \|\nabla^5z_R\|_2\le CR^{-2+s_0-5}\|v\|
                     =CR^{-1-a}\|v\|,
 \qquad
 \|\nabla^6z_R\|_2\le CR^{-2}\|v\|.
\]
In the second estimate, terms with at most five derivatives on $v$ are
actually $O(R^{-2-a})$; the sixth derivative term is $O(R^{-2})$.
Since $s_0=5a+6(1-a)$, interpolation yields the explicit bound
\begin{equation}\label{lin:fractional-tail}
 \|z_R\|_{\dot H^{s_0}}
 \le\|\nabla^5z_R\|_2^a\|\nabla^6z_R\|_2^{1-a}
 \le CR^{-2+a-a^2}\|v\|.
\end{equation}
The identical proof uses Fourier sums on a torus.  Its remaining
regularized term is bounded separately by
\begin{equation}\label{lin:zero-mode-tail}
 L_s^{-s_0}\|z_R\|_2
 \le CR^{-2}(R/L_s)^{s_0}\|v\|
 \le CR^{-2}\|v\|,
\end{equation}
with a harmless geometric constant when $R\le C_0L_s$.
Both exponents are summable over dyadic annuli; indeed
$-2+a-a^2\le-7/4$ for $0<a<1$.  Summation uses the Sobolev triangle
inequality, not a locality assertion for fractional derivatives.

It remains to estimate a fixed interior region.  Its low-order product
is controlled by a local $H^6$ norm, and the lower top-order Leibniz
terms by a local $H^{k-1}$ norm.  Choose a slightly larger fixed ball
and a cutoff equal to one on the region.  The usual interpolation
inequality on these two balls gives, for each $m<k$,
\[
 \|v\|_{H^m(B_R)}
 \le\delta\|v\|_{H^k(B_{2R})}
                 +C_{R,m,\delta}\|v\|_{L^2(B_{2R})}.
\]
For completeness, this follows from compactness of the embedding
$H^k(B_{2R})\to H^m(B_R)$: a failure after normalizing the left side
would give a bounded sequence whose $L^2$ norm tends to zero but whose
$H^m$ norm does not.  Local extension and Fourier truncation give this
compactness.  By \eqref{lin:local-derivatives}, the $H^k$ norm on the
fixed larger ball is bounded by $C_R\|v\|$.  Choose the tail radius
first and $\delta$ next.  This proves \eqref{lin:compact-bound}, including
the interior contribution of the regularized torus term.
\end{proof}

\begin{proposition}[Energy observation inequality]\label{lin:observation}
There are $c>0$, $C<\infty$, and $R_0<\infty$ such that every whole-space
linear solution, and every torus linear solution with sufficiently large
starting scale, satisfies
\begin{equation}\label{lin:energy-observation}
 \frac{d}{ds}\|v(s)\|^2
 \le-c\|v(s)\|^2+C\|v(s)\|_{L^2(B_{R_0})}^2.
\end{equation}
The derivative can equivalently be understood through the integrated
inequality.  The constants are independent of the torus scale and of
the slab length.  In particular, for $t\ge0$,
\begin{equation}\label{lin:integrated-observation}
 \|v(t)\|^2\le e^{-ct}\|v(0)\|^2
   +C\int_0^t e^{-c(t-s)}\|v(s)\|_{L^2(B_{R_0})}^2\,ds.
\end{equation}
\end{proposition}

\begin{proof}
The two squared free norm rates, on both kinds of space, are
\begin{equation}\label{lin:free-rates}
 6-s_0-2a=-a,\qquad 6-k-2a.
\end{equation}
For the first periodic norm this also uses
$L_s^{-2}=e^{-s}L^{-2}$; thus the zero mode has exactly the same first
rate.  At order $k$, the leading potential term contributes at most
$2B\|\nabla^kv\|_2^2$.  This constant does not grow with $k$, because
the same pointwise operator acts on each ordered derivative component.
All combinatorial constants occur only in the lower terms covered by
Lemma~\ref{lin:compact-error}.  The entire fractional potential term is
also covered there.  Consequently, writing $N_0,N_k$ for the two norm
components, one obtains
\[
 \frac{d}{ds}(N_0^2+N_k^2)
 \le-aN_0^2-(k+2a-6-2B)N_k^2
       +C\|v\|\big(\varepsilon\|v\|
                         +C_\varepsilon\|v\|_{L^2(B_{R_\varepsilon})}\big).
\]
Choose $\varepsilon$ after the fixed $k$, and apply Young's inequality.
Condition~\eqref{lin:k-choice} leaves a positive $c$.

This computation does not demand two additional derivatives of the
initial datum.  On a short step pull the mild solution back by its free
propagator.  Strong continuity and the bounded potential give first
variation $-i\mathcal V_*v$ in the relevant space, while the two free
squared norms acquire exactly the exponential factors in
\eqref{lin:free-rates}.  Taking the norm variation gives the displayed
inequality for arbitrary mild solutions.  Equivalently one can first
use smooth data, integrate the inequality, and pass to general data by
the uniform mild-solution bounds.  This proves
\eqref{lin:integrated-observation} as well.
\end{proof}

\subsection{Actual decay on the whole-space stable subspace}

Let $E:\mathbb R^{14}\to Y$ have columns
\begin{equation}\label{lin:E}
 iQ,\quad \partial_1Q,\ldots,\partial_{12}Q,\quad
 (a-ib+\tfrac12y\cdot\nabla)Q,
 \qquad
 \Lambda=\operatorname{diag}(0,\tfrac12,\ldots,\tfrac12,1).
\end{equation}
These columns are independent symmetry vectors, as established in
Theorem~\ref{spectrum:classification}.

The passage from dissipation modulo a compact term to stable semigroup
decay also appears in \cite[Section~4.2]{BuckEtAl2026}. Here the full
mode classification identifies the complementary spectral space with
the fourteen symmetry directions above.

\begin{theorem}[Stable whole-space evolution]\label{linear:decay}
There is a bounded real-linear map $\pi:Y\to\mathbb R^{14}$ such that
\begin{equation}\label{lin:spectral-coordinates}
 \pi E=I,\qquad S(s)E=Ee^{s\Lambda},\qquad
 \pi S(s)=e^{s\Lambda}\pi\quad(s\in\mathbb R).
\end{equation}
For some $C_S<\infty$ and $\eta>0$,
\begin{equation}\label{lin:stable-decay}
 \|S(s)v\|_Y\le C_Se^{-\eta s}\|v\|_Y,
 \qquad s\ge0,\quad \pi v=0.
\end{equation}
Thus $E\pi$ is the symmetry spectral projection.  Constants depend on
the fixed profile, $a$, and $k$.
\end{theorem}

\begin{proof}
We first obtain a statement about the bounded time-one operator
$T=S(1)$, rather than deducing decay from a generator spectral gap.
If $v_j\rightharpoonup0$ in $Y$ and $\|v_j\|\le1$, then for every fixed
$s$ the functions $S(s)v_j$ converge strongly to zero in
$L^2(B_{R_0})$.  Indeed they converge weakly in $Y$, and the restriction
from $Y$ to that $L^2$ space is compact by
\eqref{lin:local-derivatives} and local Sobolev compactness.  Slab bounds
and dominated convergence in \eqref{lin:integrated-observation} give
\begin{equation}\label{lin:essential-contraction}
 \limsup_j\|Tv_j\|\le e^{-c/2}.
\end{equation}
The same conclusion holds on the complexification, by adding the two
real energy inequalities.

Let $F_N$ be increasing finite-dimensional orthogonal projections
converging strongly to the identity.  Formula~\eqref{lin:essential-contraction}
implies $\limsup_N\|T(I-F_N)\|\le e^{-c/2}$: otherwise unit vectors
in the complements of $F_N$ would give a weakly null counterexample.
For some $N$ we therefore have
\[
 T=B_0+K_0,\qquad \|B_0\|=q<1,\qquad \operatorname{rank}K_0<\infty.
\]
Outside $|z|\le q$, invert $z-B_0$ and write
$K_0=UV$ through a finite-dimensional space.  The resolvent of $T$
is determined by the finite matrix
$I-V(z-B_0)^{-1}U$.  Its determinant is analytic and tends to one as
$|z|\to\infty$.  Thus the spectrum there consists of isolated
eigenvalues with finite-dimensional generalized eigenspaces; there
are only finitely many outside any larger disk.  This also proves the
usual spectral decomposition in this region by the finite-matrix
resolvent formula.

Choose $r$ with
\[
 \max\{q,e^{-1/32}\}<r<1
\]
and with $|z|=r$ disjoint from the spectrum.  Every spectral subspace
of $T$ outside this circle is preserved by $S(s)$, because $S(s)$
commutes with the time-one resolvent and its contour integrals.
The restricted continuous group on a finite-dimensional space is a
matrix exponential.  Each of its generator eigenvalues $\lambda$
satisfies $\operatorname{Re}\lambda=\log|z|>\log r>-1/32$ for the corresponding
time-one eigenvalue $z$.

Its generator equations are the distributional equations for
\eqref{lin:potential}, as follows by differentiating
\eqref{lin:mild} against compactly supported smooth tests.  Local
elliptic regularity gives smoothness of eigenvectors and chain
vectors.  Projection onto any spherical harmonic gives regular
radial solutions of the mode equation in
Theorem~\ref{spectrum:classification}.  The unweighted radial $k$th
derivative is in $L^2(r^{11}dr)$: along each ray $\partial_r^k$ is a
contraction of the ordered tensor $\nabla^k$, and angular projection
is bounded on $L^2$ of each sphere.  Smoothness of the projection at
the origin follows by projecting the Taylor polynomials, using the
harmonic decomposition of homogeneous polynomials.  The
top-derivative criterion in Theorem~\ref{spectrum:classification}
excludes the fast exterior branches, including in a generalized
eigenvector after subtraction of the differentiated slow basis.
That theorem therefore leaves precisely the symmetry vectors
\eqref{lin:E} and excludes all nontrivial Jordan chains.  Completeness
of spherical harmonics gives the statement for the original vectors;
a nontrivial chain would have a nonzero harmonic projection of its
leading vector and hence would contradict the same classification.

Conversely all columns of $E$ lie in $Y$ by
\eqref{lin:symbol-norm}; the time column even has symbol order
$-2a-2$, since its leading slow power cancels.  They are generator
domain vectors with the asserted eigenvalues.  To check the domain
claim directly, multiply a column by a smooth cutoff at radius $R$.
Both the discarded symbol tail and the commutator
\[
 i[\Delta,\chi(y/R)]E-\tfrac12(y\cdot\nabla\chi(y/R))E
\]
tend to zero in $Y$ by the same dyadic estimates.  The smooth compactly
supported approximations thus converge in the generator graph norm.
It follows that the entire spectral subspace outside $|z|=r$ is
$\operatorname{ran}E$, with the action in
\eqref{lin:spectral-coordinates}.  Complex conjugation in the
complexification preserves its projection, so this projection is the
complexification of a bounded real projection $E\pi$.

On its complement, the spectrum of $T$ lies in $|z|\le r$.  Choose
$r<r_1<1$.  The contour formula for powers on $|z|=r_1$ gives
\[
 \|T^n|_{\ker\pi}\|
 \le r_1^{n+1}\max_{|z|=r_1}
                  \|(z-T|_{\ker\pi})^{-1}\|.
\]
Bounded propagation for the remaining time in $[0,1]$ proves
\eqref{lin:stable-decay}, with $\eta=-\log r_1>0$ and an enlarged
$C_S$.  This derivation supplies actual semigroup decay, not only a
description of the generator spectrum.
\end{proof}

\subsection{A transfer proposition for expanding tori}

Define the cellwise cutoff symmetry map $E_L=\chi(y/L)E$.
The bounds already proved imply
\begin{equation}\label{lin:EL-convergence}
 \sup_{L\ge1}\|E_L\|<\infty,\qquad J_LE_L\longrightarrow E
       \quad\hbox{in }\mathcal L(\mathbb R^{14},Y).
\end{equation}
Indeed $J_LE_L=\chi(y/L)^2E$, and its omitted dyadic tail tends to
zero.  For all sufficiently large $L$ put
\begin{equation}\label{lin:piL}
 C_L=\pi J_LE_L,\qquad
 \pi_L=C_L^{-1}\pi J_L,\qquad P_L=E_L\pi_L.
\end{equation}
Then $C_L\to I$, $\pi_LE_L=I$, and $P_L$ is a projection.  The
operator norms of $E_L,\pi_L,P_L$, and $I-P_L$ are uniformly bounded.
Let $\mathcal F_L^{(M)}:Y_L\to Y_{Le^{M/2}}$ denote evolution by
\eqref{lin:torus-evolution} from time zero to time $M$.

\begin{proposition}[Uniform expanding-torus transfer]\label{linear:torus}
For each fixed $M>0$ the step maps $\mathcal F_L^{(M)}$ are uniformly
bounded for all sufficiently large $L$.  Writing $L'=Le^{M/2}$,
\begin{equation}\label{lin:coordinate-defect}
 \lim_{L\to\infty}
 \|\pi_{L'}\mathcal F_L^{(M)}-e^{M\Lambda}\pi_L\|
 _{\mathcal L(Y_L,\mathbb R^{14})}=0.
\end{equation}
There exists $M_0<\infty$ such that for every fixed $M\ge M_0$
there is $L_0(M)<\infty$ for which
\begin{equation}\label{lin:torus-stable-contraction}
 \|(I-P_{L'})\mathcal F_L^{(M)}|_{\ker\pi_L}\|
       _{\mathcal L(Y_L,Y_{L'})}\le\frac18,
 \qquad L\ge L_0(M).
\end{equation}
Moreover $\|(I-P_{L'})\mathcal F_L^{(M)}E_L\|\le C_M$.
Any prescribed positive bound on the coordinate defect can be imposed
by enlarging $L_0(M)$ after $M$ has been fixed.  These bounds persist
after restarting a step at any larger scale.

After identifying torus spaces by dilation, the maps above and the
step propagator are continuous when the scale and the input vary
continuously.  This last assertion is strong continuity for varying
inputs, and does not assert operator-norm continuity of Schr\"odinger
time evolution.
\end{proposition}

\begin{proof}
We give the compactness and uniqueness argument in detail.  Fix $M$
and a sequence $L_j\to\infty$ with initial data $v_j(0)$ bounded in
$Y_{L_j}$.  Write $L_{j,s}=L_je^{s/2}$.  Uniform slab bounds and
\eqref{lin:local-derivatives} give boundedness in local $H^k$, uniformly
for $0\le s\le M$.  On each fixed ball the equation also bounds
$\partial_sv_j$ in $H^{k-2}$; the drift coefficient is bounded on
that ball.  Local compactness, interpolation, and Arzel\`a--Ascoli
therefore give a subsequence converging in
$C([0,M];L^2(B_R))$ for every fixed $R$, to a function $v$.
One may obtain the stronger convergence in local $H^{k-1}$, but it
is unnecessary here.

The localized functions $J_{L_{j,s}}v_j(s)$ are bounded in $Y$ by
\eqref{lin:J-bound}.  At each $s$, every weakly convergent subsequence
has the same distributional limit $v(s)$, since $\chi(y/L_{j,s})=1$
on each fixed ball for all large $j$.  Thus
\begin{equation}\label{lin:weak-local-limit}
 J_{L_{j,s}}v_j(s)\rightharpoonup v(s)\quad\hbox{in }Y
 \quad(0\le s\le M),\qquad
 \|v(0)\|_Y\le C_J\liminf_j\|v_j(0)\|_{Y_{L_j}}.
\end{equation}
The limit is weakly continuous in $Y$: boundedness and its local
continuity identify the weak limit along every converging time
sequence.  As $Y$ is separable, it is strongly measurable and bounded
on this slab.  The torus coefficients agree with the whole-space
coefficients on each fixed ball eventually, so local passage in the
equation gives $\partial_sv=\mathcal Av$ in distributions.

This distributional solution is the whole-space mild solution.
To see this without a propagation assertion, test against Schwartz
functions transported by the free adjoint.  Formula~\eqref{lin:free-flow}
preserves Schwartz functions, uniformly on compact time intervals.
The global $L^\infty$ bound supplied by $Y$, and the at-most-linear
growth of the drift, justify removal of spatial cutoffs in these
tests.  Integration in time gives
\[
 S_0(-s)v(s)=v(0)-i\int_0^s S_0(-t)\mathcal Vv(t)\,dt
\]
as tempered distributions.  The integral is a Bochner integral in
$Y$ by boundedness of $\mathcal V$ and strong measurability.  The
distributional identity is therefore an identity in $Y$ and is
equivalent to \eqref{lin:mild}.  Uniqueness gives
\begin{equation}\label{lin:limit-identification}
 v(s)=S(s)v(0).
\end{equation}

Finite dimensionality, \eqref{lin:weak-local-limit}, and $C_L\to I$
now imply convergence of the initial and final coordinates to
$\pi v(0)$ and $\pi S(M)v(0)$, respectively.  Their difference in
\eqref{lin:coordinate-defect} tends to zero by
\eqref{lin:spectral-coordinates}.  This proves convergence in
operator norm: failure would furnish a sequence of unit initial
vectors witnessing a fixed positive lower bound, to which exactly
the preceding subsequence argument applies.

For the contraction, take such a sequence with
$\|v_j(0)\|_{Y_{L_j}}=1$ and $\pi_{L_j}v_j(0)=0$.  The limit satisfies
$\pi v(0)=0$, so \eqref{lin:stable-decay} gives
\[
 \|v(s)\|_Y\le C_SC_Je^{-\eta s}.
\]
The constants in this estimate are independent of the fixed slab
length $M$.  Apply the integrated observation inequality
\eqref{lin:integrated-observation} \emph{before} passing to a local
limit.  The full initial norm is still one in its first term.  Local
uniform convergence passes only the observation integral to the
limit, giving
\begin{equation}\label{lin:escaped-norm}
 \limsup_j\|\mathcal F_{L_j}^{(M)}v_j(0)\|_{Y_{L_je^{M/2}}}^2
 \le D(M):=e^{-cM}+C_1\int_0^M
                         e^{-c(M-s)}e^{-2\eta s}\,ds.
\end{equation}
Here $C_1$ uses only the fixed observation embedding, $C_J$, $C_S$,
and the energy constant; it does not contain the slab bound $C_M$.
Any norm that escaped to spatial infinity remains charged in the
first term $e^{-cM}$.  In particular a zero local limit causes no loss
in this estimate.

The quantity $D(M)$ tends to zero as $M\to\infty$; when $c=2\eta$
its integral is $Me^{-cM}$, and otherwise it is
$(e^{-2\eta M}-e^{-cM})/(c-2\eta)$.  Let $C_P$ uniformly bound
$\|I-P_L\|$ for sufficiently large $L$.  Choose $M_0$ so that
$C_P\sqrt{D(M)}<1/16$ for every $M\ge M_0$.  For each such fixed $M$,
a failure of \eqref{lin:torus-stable-contraction} along arbitrarily
large $L$ would supply unit stable vectors contradicting
\eqref{lin:escaped-norm}.  This proves the required order of choices:
first $M$, then $L$.  The stated bound on the map from $E_L$ follows
from the uniform step and projection bounds.  Restarting a step uses
the same formulas with its current scale as the initial scale.

Finally identify $v\in Y_L$ with $v(Lx)$ on $\mathbb T^{12}$.
On a compact range of positive $L$, these norms are equivalent to
$H^k$, continuously in $L$.  The free step on this fixed torus has
Fourier phases
$\exp(-i|n|^2L^{-2}(1-e^{-s}))$ and scalar factor $e^{(-a+ib)s}$.
Dominated convergence of the weighted Fourier sum proves strong
continuity for continuously varying inputs.  The cutoff coefficients
are smooth and continuous in $L$ on each finite slab; successive
approximation in the integral equation transfers this continuity to
the full propagator.  Dilation is strongly continuous on $Y$, which
gives the same assertion for $J_L,E_L$ and then for $\pi_L,P_L$ by
\eqref{lin:piL}.  This completes the proof.
\end{proof}
